\documentclass[11pt,a4paper]{article}

\usepackage[T1]{fontenc}
\usepackage[utf8]{inputenc}
\usepackage[brazil]{babel}
\usepackage{lmodern}
\usepackage{microtype}
\usepackage[margin=2cm]{geometry}
\usepackage{hyperref}
\usepackage{titlesec}
\usepackage{xcolor}

\hypersetup{
  colorlinks=true,
  linkcolor=black,
  urlcolor=black,
  citecolor=black,
  pdfborder={0 0 0}
}

\definecolor{sectioncolor}{HTML}{111111}
\definecolor{accentcolor}{HTML}{1A1A1A}

\pagestyle{empty}
\setlength{\parindent}{0pt}
\setlength{\parskip}{5pt}
\linespread{1.05}
\selectfont

\titleformat{\section}
  {\bfseries\large\color{sectioncolor}}
  {}{0pt}{}[\vspace{-2pt}{\color{sectioncolor}\titlerule[0.8pt]}]
\titlespacing*{\section}{0pt}{12pt}{8pt}

\newcommand{\cvrole}[4]{%
  \vspace{4pt}%
  {\bfseries\color{accentcolor}#1}\hfill{\itshape #2}\\[-2pt]%
  {\itshape #3}\hfill #4\\[6pt]%
}

\begin{document}

\begin{center}
  {\fontsize{20}{24}\selectfont\bfseries Kamille Hartmann Maccagnan}\\[8pt]
  {\large Analista de Testes \textbar{} Validação de Sistemas \textbar{} QA}\\[10pt]
  Campo Comprido, Curitiba-PR \enspace\textbullet\enspace (41) 9 9928-4424 \enspace\textbullet\enspace
  \href{mailto:kamillehartmannm@gmail.com}{kamillehartmannm@gmail.com}\\[3pt]
  \href{https://linkedin.com/in/kamille-hartmann-maccagnan/}{linkedin.com/in/kamille-hartmann-maccagnan/}
\end{center}

\vspace{6pt}

\section{Resumo Profissional}

Profissional com experiência em testes manuais de sistemas, validação pré-implantação e documentação de processos em ambiente de varejo. Atua na execução de testes funcionais antes da liberação de sistemas para a rede de lojas, evidenciando resultados, reportando incidentes e elaborando POPs para apoio à operação. Vivência com ServiceNow para registro e acompanhamento de chamados, validação de dados e comunicação clara de problemas para equipes técnicas. Perfil analítico, atento a detalhes, comunicativo e colaborativo, com proatividade e interesse em evoluir na área de qualidade de software.

\section{Experiência Profissional}

\cvrole{Anjuss}{Analista de Suporte de TI}{07/2025 -- Atual}

Atuo na realização de testes manuais do sistema antes da liberação para a rede de lojas, validando funcionalidades, identificando falhas e reportando incidentes para a equipe de T.I. com comunicação clara e objetiva. Elaboro e atualizo documentações e POPs do sistema para as lojas e para a equipe técnica, evidenciando procedimentos e resultados esperados das rotinas validadas. Lidero treinamentos de novos funcionários após a homologação das funcionalidades, verificando a correta aplicação dos processos em ambiente operacional. Apoio o monitoramento de indicadores em Power BI, com validação de dados e identificação de inconsistências que possam impactar a operação das lojas.

\cvrole{HORSE do Brasil}{Estagiária de TI (IS/IT LATAM)}{07/2024 -- 07/2025}

Atuei no suporte e acompanhamento de demandas em ambiente multinacional, utilizando o ServiceNow para registro, acompanhamento e reporte de incidentes, apoiando a organização de chamados e o alinhamento entre usuários e áreas técnicas. Participei de projetos de governança de dados e melhoria contínua, com documentação de processos, verificação de informações e comunicação de ocorrências para resolução. Elaborei relatórios e dashboards em Power BI para acompanhamento de indicadores, com atenção à qualidade e consistência dos dados. Atuei como ponte entre áreas de negócio, operações e stakeholders internacionais, colaborando com equipes multidisciplinares.

\cvrole{Restaurante Familiar}{Analista de Dados e Suporte Técnico}{01/2023 -- 07/2024}

Atuei no suporte técnico e na validação de informações operacionais e financeiras, com verificação de consistência de dados e identificação de desvios em relatórios e indicadores. Estruturei bases de dados e dashboards em Power BI para acompanhamento de resultados, aplicando regras de validação para garantir confiabilidade das análises. Apoiei a gestão na identificação de problemas recorrentes em processos operacionais e na comunicação de necessidades de correção entre áreas operacionais e administrativas.

\section{Competências Técnicas}

\textbf{Testes e Qualidade:} Testes manuais de software, validação pré-implantação, testes funcionais, evidência de testes, reporte de incidentes, atenção a detalhes, documentação de processos, POPs, melhoria contínua.

\textbf{Ferramentas e Sistemas:} ServiceNow, Power BI, Excel, SQL, Power Automate, Google Sheets, Pacote Office, Trello, Monday.

\textbf{Análise e Dados:} Verificação de consistência de dados, indicadores operacionais, relatórios gerenciais, capacidade analítica, raciocínio lógico.

\textbf{Colaboração:} Trabalho em equipe, interface entre usuários e áreas técnicas, treinamento e capacitação, comunicação com stakeholders, gestão de demandas.

\textbf{Metodologias:} Scrum, Kanban.

\textbf{Idiomas:} Inglês Avançado \enspace\textbullet\enspace Espanhol Básico \enspace\textbullet\enspace Coreano Básico

\section{Projetos e Conquistas}

1º lugar no Transformation Day Renault 2023 com solução orientada a dados e melhoria de processos. Participação no Hackathon Bosch 2023 com foco em resolução de problemas. Execução de testes manuais e validação de sistemas antes da implantação na rede de lojas na Anjuss. Desenvolvimento de documentações e POPs para apoio à operação e evidência de rotinas validadas.

\section{Formação Acadêmica}

\textbf{PUCPR} \hfill Curitiba, PR\\
Graduação em Gestão da Tecnologia da Informação \hfill Previsão de conclusão: 06/2026\\[6pt]

\textbf{Escola Conquer} \hfill Curitiba, PR\\
Pós-graduação em Liderança e Gestão em Tecnologia \hfill Conclusão: 2024\\[6pt]

\textbf{Universidade Tuiuti do Paraná} \hfill Curitiba, PR\\
Graduação em Medicina Veterinária \hfill 06/2019\\[6pt]

\textbf{Qualittas} \hfill Curitiba, PR\\
Pós-graduação em Clínica Médica e Cirúrgica de Animais Exóticos e Pets Não Convencionais \hfill 12/2022

\end{document}
